% !TeX root = ../main.tex

\chapter{Meccanismi di Reliability e Scalabilità}
\label{cap:reliability_scalability}

Il presente capitolo illustra le soluzioni applicate per gestire i temporanei disservizi e per scalare le risorse al bisogno. L'architettura fa ampio uso dei comportamenti \emph{out-of-the-box} forniti dalle librerie di base, arricchiti da pattern architetturali pensati per confinare gli errori e rendere il sistema reattivo al traffico.

% ---------------------------------------------------------
\section{Resilienza di Rete}
\label{sec:resilienza_rete}

Nei sistemi distribuiti basati su messaggistica (come i client implementati in \texttt{drone\_simulator.py}), la sconnessione temporanea del broker MQTT è un'evenienza da gestire nativamente.

In fase di avvio (funzione \texttt{run()} del simulatore), i droni eseguono un ciclo di retry custom: qualora il broker non sia ancora schedulato (ad esempio durante un deployment asincrono in un cluster Kubernetes), la connessione fallisce e viene avviato un poller ritardato di 5 secondi tramite \texttt{time.sleep}, ripetuto fino al successo dell'handshake.

A regime, la maggior parte della resilienza di rete è demandata ai meccanismi interni della libreria \texttt{paho-mqtt} e al suo metodo \texttt{loop\_start()}. In caso di caduta momentanea del servizio Mosquitto, il client Python intercetta il blocco della socket e avvia autonomamente una strategia di \textit{exponential backoff} per i tentativi di riconnessione, ripristinando la sottoscrizione ai topic operativi (\texttt{comandi/\#}) senza introdurre logiche custom che rischierebbero di sovraccaricare il thread principale.

% ---------------------------------------------------------
\section{Graceful Degradation del Server Centrale}
\label{sec:graceful_degradation}

Nel sistema, \texttt{central\_server.py} gestisce sia la persistenza delle telemetrie su InfluxDB sia l'esposizione della dashboard Flask in tempo reale per la supervisione della flotta. Per evitare che un downtime del database causi il collasso del monitoraggio live, è stato adottato un approccio ibrido con degradazione controllata.

Il pattern prevede un \textit{dual-write} all'arrivo di ogni messaggio MQTT (handler \texttt{on\_message}), strutturato secondo la seguente logica:

\begin{lstlisting}[language=Python, caption={Schema logico del meccanismo di Graceful Degradation nel server centrale.}, label={lst:graceful_degradation}]
def on_message(client, userdata, msg):
    # 1. Aggiornamento immediato dello stato in-memory (per la dashboard live)
    state["drones"][drone_id] = parsing(msg.payload)

    # 2. Scrittura persistente su InfluxDB isolata da eccezioni
    try:
        write_api.write(bucket, org, record)
    except InfluxDBError as e:
        # L'eccezione viene neutralizzata per preservare i servizi realtime
        log_warning("InfluxDB non raggiungibile. Degradazione del servizio log storici.")
\end{lstlisting}

Se il container InfluxDB risulta non accessibile, l'eccezione viene neutralizzata nel blocco \texttt{try/except} interno all'handler. Di conseguenza, il sistema perde il salvataggio dei log storici (degradazione del servizio) ma mantiene perfettamente operativi il tracciamento real-time dei pacchetti e i cruscotti per gli operatori (comportamento \emph{graceful}), isolando il guasto al solo storage layer senza bloccare la coda delle missioni.

% ---------------------------------------------------------
\section{Scalabilità Dinamica e Safety Limits Kubernetes}
\label{sec:scalabilita_dinamica}

L'architettura supera lo scaling passivo (basato unicamente su carico CPU/RAM) delegando la valutazione dei bisogni logistici agli agenti AI dotati di capacità MCP.

In particolare, il \texttt{LogisticAgent} e l'\texttt{HealthAgent} dispongono di un tool per richiedere dinamicamente più droni e smaltire la coda (\texttt{pending\_orders}). Internamente, questo tool richiama la funzione \texttt{scale\_drone\_deployment} in \texttt{drone\_mcp\_layer.py}, che interagisce via API HTTP con il control-plane di Kubernetes. Sfruttando la libreria client ufficiale per Python, la richiesta genera un'operazione di tipo \texttt{patch\_namespaced\_deployment\_scale} sul Deployment \texttt{drone-simulator}, variando in tempo reale il numero di repliche a disposizione in base alla coda dei pacchi.

A tutela dell'infrastruttura (\textit{Operational Safety}), lo scaling dinamico è sottoposto a policy di sicurezza:

\begin{lstlisting}[language=Python, caption={Guardrail di sicurezza applicato alle richieste di scaling dell'agente.}, label={lst:safety_limits_scale}]
# Verifica dei vincoli architetturali prima del patching su Kubernetes
if requested_replicas > POLICY_LIMITS:
    # L'azione viene bloccata ed escalata all'operatore umano
    trigger_human_escalation(requested_replicas)
else:
    # Esecuzione della patch sul Deployment
    k8s_client.patch_namespaced_deployment_scale(name="drone-simulator", ...)
\end{lstlisting}

Se l'agente tenta un'espansione del numero di repliche oltre la soglia prestabilita \texttt{POLICY\_LIMITS} (vedi Capitolo~\ref{cap:safety_guardrails}), l'azione viene intercettata, bloccata ed escalata al \textit{Human Approval Manager}, evitando consumi cloud incontrollati dovuti a potenziali allucinazioni o comportamenti imprevisti del modello.
